% Nexus -- Informe Final (Stage III)
% Compile with:  pdflatex Stage-III-Report.tex  (run twice for the ToC)
% Requires a TeX distribution with: listings, xcolor, booktabs, geometry,
% hyperref, lmodern and (optionally) babel-spanish.
\documentclass[11pt,a4paper]{article}

\usepackage[T1]{fontenc}
\usepackage[utf8]{inputenc}
\usepackage{lmodern}
\usepackage[spanish,es-noshorthands]{babel}
\usepackage[a4paper,margin=2.5cm]{geometry}
\usepackage{xcolor}
\usepackage{listings}
\usepackage{booktabs}
\usepackage{longtable}
\usepackage{enumitem}
\usepackage{hyperref}
\hypersetup{colorlinks=true,linkcolor=blue!50!black,urlcolor=blue!50!black}

% ---- Code listing style (mirrors the Stage I document) --------------------
\definecolor{nexusbg}{rgb}{0.97,0.97,0.97}
\definecolor{nexuskw}{rgb}{0.00,0.20,0.60}
\definecolor{nexuscm}{rgb}{0.00,0.50,0.00}

\lstdefinelanguage{Nexus}{
  morekeywords={graph,kind,directed,undirected,traits,weighted,capacitated,
    nodes,edges,groups,constraints,assert,connected,strongly_connected,
    acyclic,reachable,tree,binary_tree,rooted_at,forall,in,derive,from,using,
    transpose,induced_subgraph,remove_self_loops,underlying,analysis,on,run,
    as,export,result,to,dot,json,root,source,sink,terminal,degree,indegree,
    outdegree,shortest_path,topological_sort,components,scc,mst,max_flow,
    weight,capacity},
  sensitive=true,
  morecomment=[l]{//},
}

\lstset{
  basicstyle=\ttfamily\small,
  backgroundcolor=\color{nexusbg},
  keywordstyle=\color{nexuskw}\bfseries,
  commentstyle=\color{nexuscm}\itshape,
  frame=single, framerule=0pt, framesep=6pt,
  breaklines=true, columns=fullflexible, showstringspaces=false,
  xleftmargin=4pt, xrightmargin=4pt, aboveskip=1em, belowskip=1em,
  language=Nexus,
}

\newcommand{\code}[1]{\texttt{#1}}

\begin{document}

% --------------------------------------------------------------------------
\begin{center}
  {\Huge\bfseries Nexus}\\[0.4em]
  {\Large Diseño e Implementación de un Lenguaje}\\[0.8em]
  {\large Informe Final --- Stage III}\\[0.6em]
  Autómatas, Teoría de Lenguajes y Compiladores\\
  Ingeniería Informática --- ITBA
\end{center}
\vspace{1em}

\begin{center}
\begin{tabular}{@{}llll@{}}
\toprule
\textbf{Nombre} & \textbf{Apellido} & \textbf{Legajo} & \textbf{E-mail} \\
\midrule
Pablo   & Gorostiaga & 65148 & pgorostiaga@itba.edu.ar \\
Franco  & Pampuri    & 65552 & fpampuri@itba.edu.ar \\
Jonas   & Glaubart   & 65790 & jglaubart@itba.edu.ar \\
Eugenio & Migliaro   & 65508 & emigliaro@itba.edu.ar \\
\bottomrule
\end{tabular}
\end{center}
\vspace{1.5em}

\tableofcontents
\newpage

% ==========================================================================
\section{Introducción}

\emph{Nexus} es un compilador para un DSL (\emph{Domain-Specific Language})
orientado a grafos, desarrollado para la materia Autómatas, Teoría de Lenguajes
y Compiladores. El proyecto se entregó en tres etapas que se reflejan en este
informe:

\begin{enumerate}[noitemsep]
  \item \textbf{Stage I --- Diseño:} concepción del dominio, la gramática y la
        sintaxis del lenguaje.
  \item \textbf{Stage II --- Frontend:} análisis léxico (Flex), análisis
        sintáctico (Bison) y construcción del AST.
  \item \textbf{Stage III --- Backend:} análisis semántico, generación de
        código y la documentación final.
\end{enumerate}

El objetivo de Stage III fue completar el \emph{pipeline} del compilador:
validar el programa de entrada contra las reglas del dominio (análisis
semántico), emitir un artefacto ejecutable (generación de código) y documentar
las decisiones de diseño e implementación tomadas a lo largo de todo el
proyecto.

El compilador acepta programas Nexus que declaran grafos con nodos, aristas,
grupos y restricciones; definen derivaciones (transformaciones de grafos); y
especifican análisis (algoritmos de grafos con exportación de resultados). La
salida es un programa Python que, junto con la biblioteca de \emph{runtime}
\code{nexus\_runtime.py}, construye los grafos, valida las restricciones,
ejecuta los algoritmos y exporta los resultados a formato DOT o JSON.

% ==========================================================================
\section{Modelo Computacional}

\subsection{Dominio}

El dominio de Nexus es el \textbf{modelado, validación, transformación y
análisis de grafos}. El lenguaje describe grafos como objetos matemáticos de
manera declarativa, indicando su estructura, sus propiedades y los análisis
que se desean ejecutar sobre ellos. En concreto, Nexus permite:

\begin{itemize}[noitemsep]
  \item \textbf{Declarar grafos} dirigidos o no dirigidos, con \emph{traits}
        opcionales (\code{weighted}, \code{capacitated}).
  \item \textbf{Definir estructura:} nodos con atributos semánticos opcionales
        (\code{root}, \code{source}, \code{sink}, \code{terminal}), aristas con
        peso y/o capacidad, y grupos de nodos.
  \item \textbf{Especificar restricciones} sobre la estructura del grafo:
        conectividad, aciclicidad, alcanzabilidad, estructura de árbol y
        predicados cuantitativos sobre grado.
  \item \textbf{Derivar nuevos grafos} mediante transformaciones: transposición,
        subgrafo inducido, eliminación de \emph{self-loops} y grafo subyacente.
  \item \textbf{Ejecutar algoritmos:} camino mínimo, orden topológico,
        componentes conexas, componentes fuertemente conexas, árbol de
        expansión mínima y flujo máximo.
  \item \textbf{Exportar} grafos y resultados a DOT (para visualización con
        Graphviz) o JSON.
\end{itemize}

El foco del diseño estuvo puesto en un núcleo reducido pero expresivo, con una
especificación clara, consistente y formalmente verificable: cada construcción
del lenguaje tiene precondiciones de dominio que el compilador puede comprobar
estáticamente.

\subsection{Lenguaje}

\subsubsection{Gramática}

El lenguaje queda definido por una gramática
$G = \langle \Sigma, N, \Pi, S \rangle$ de tipo 2 (libre de contexto). El
analizador léxico (Flex) define el alfabeto $\Sigma$ a partir de los lexemas
del lenguaje, mientras que el analizador sintáctico (Bison) define el conjunto
de no-terminales $N$, las producciones $\Pi$ y el símbolo inicial $S$. Ambas
fases en conjunto conforman la gramática que genera el lenguaje. La gramática
es LALR(1) y libre de conflictos.

El símbolo inicial $S$ deriva un \code{Program}: una o más declaraciones de
nivel superior, cada una de las cuales es una declaración de grafo
(\code{graph}), una derivación (\code{derive}) o un bloque de análisis
(\code{analysis}).

\subsubsection{Sintaxis general}

\begin{lstlisting}
graph Build:
    kind directed
    traits weighted
    nodes:
        lexer
        parser
        ast
        emit
    edges:
        lexer -> parser weight 2
        parser -> ast weight 3
        ast -> emit weight 4
    groups:
        core = {lexer, parser, ast}
        sinks = {emit}
    constraints:
        assert acyclic
        assert reachable lexer -> emit
        assert forall n in sinks: outdegree(n) = 0

derive CoreBuild from Build using induced_subgraph(core)

analysis Main on Build:
    run shortest_path from lexer to emit as critical_path
    export graph to dot
    export result critical_path to json
\end{lstlisting}

Los bloques \code{traits}, \code{groups} y \code{constraints} son opcionales;
cuando aparecen deben respetar exactamente la sintaxis del lenguaje y ubicarse
dentro de la declaración del grafo correspondiente.

\subsubsection{Construcciones}

\begin{longtable}{@{}p{0.46\linewidth}p{0.50\linewidth}@{}}
\toprule
\textbf{Construcción} & \textbf{Significado} \\
\midrule
\endhead
\code{graph <id>:} & Declara un grafo con identificador único. \\
\code{kind directed|undirected} & Tipo del grafo (obligatorio). \\
\code{traits weighted capacitated} & Propiedades opcionales del grafo. \\
\code{nodes:} & Sección de declaración de nodos. \\
\code{<id> [root|source|sink|terminal]} & Nodo con atributo semántico opcional. \\
\code{edges:} & Sección de declaración de aristas. \\
\code{<u> -> <v> [weight N] [capacity N]} & Arista dirigida (peso/capacidad opcional). \\
\code{<u> -- <v> [weight N] [capacity N]} & Arista no dirigida (peso/capacidad opcional). \\
\code{groups:} & Sección de declaración de grupos. \\
\code{<id> = \{<n1>, <n2>, ...\}} & Grupo nombrado de nodos. \\
\code{constraints:} & Sección de restricciones. \\
\code{assert connected} & Grafo conexo (solo \code{undirected}). \\
\code{assert strongly\_connected} & Fuertemente conexo (solo \code{directed}). \\
\code{assert acyclic} & Grafo acíclico. \\
\code{assert reachable <u> -> <v>} & Existe un camino de $u$ a $v$. \\
\code{assert tree rooted\_at <r>} & Es árbol con raíz $r$ (solo \code{directed}). \\
\code{assert binary\_tree rooted\_at <r>} & Es árbol binario con raíz $r$ (solo \code{directed}). \\
\code{assert forall <v> in <g>: <pred>} & Predicado de grado para todo nodo del grupo. \\
\code{derive <id> from <src> using <t>} & Crea un grafo derivado. \\
\code{analysis <id> on <graph>:} & Bloque de análisis sobre un grafo. \\
\code{run <algo> [from <u> to <v>] as <r>} & Ejecuta un algoritmo y nombra su resultado. \\
\code{export graph to dot|json} & Exporta el grafo. \\
\code{export result <id> to dot|json} & Exporta un resultado. \\
\bottomrule
\end{longtable}

\noindent\textbf{Predicados de grado:} \code{degree(n)}, \code{indegree(n)},
\code{outdegree(n)} comparados con \code{=}, \code{!=}, \code{>=}, \code{<=},
\code{<}, \code{>} y un valor entero.\\
\textbf{Transformaciones:} \code{transpose}, \code{induced\_subgraph(<group>)},
\code{remove\_self\_loops}, \code{underlying}.\\
\textbf{Algoritmos:} \code{shortest\_path}, \code{topological\_sort},
\code{components}, \code{scc}, \code{mst}, \code{max\_flow}.

\subsubsection{Semántica de los atributos de nodo}

Los atributos de nodo no son meras etiquetas descriptivas, sino anotaciones con
semántica verificable por el compilador:

\begin{itemize}[noitemsep]
  \item \code{root}: vértice raíz del grafo. A lo sumo uno por grafo; en grafos
        dirigidos debe cumplir \code{indegree(root) = 0}. Toda restricción de
        \code{tree}/\code{binary\_tree} debe enraizarse en este nodo.
  \item \code{source}: origen principal de un grafo dirigido. A lo sumo uno;
        debe cumplir \code{indegree = 0} y \code{outdegree > 0}.
  \item \code{sink}: destino principal de un grafo dirigido. A lo sumo uno;
        debe cumplir \code{outdegree = 0} e \code{indegree > 0}.
  \item \code{terminal}: nodo hoja. En grafos dirigidos
        \code{outdegree = 0}; en no dirigidos \code{degree = 1}. Puede haber
        varios.
\end{itemize}

\subsubsection{Reglas estructurales}

Se permiten \emph{self-loops} pero no aristas paralelas: entre dos vértices
puede existir a lo sumo una arista (pares ordenados en grafos dirigidos, pares
no ordenados en no dirigidos). Si un grafo declara \code{assert acyclic},
el compilador rechaza todo \emph{self-loop} y todo ciclo, ya que constituyen
violaciones de la aciclicidad.

% ==========================================================================
\section{Implementación}

El \code{EntryPoint} ejecuta el \emph{pipeline} en tres fases:
\textbf{parse $\rightarrow$ análisis semántico $\rightarrow$ generación de
código}. La fase semántica sólo habilita la generación si el programa es
válido respecto del dominio.

\subsection{Frontend}

El frontend fue construido en Stage II y no requirió modificaciones de fondo en
Stage III.

\textbf{Análisis léxico (Flex).} El \emph{lexer} (\code{FlexPatterns.l})
reconoce las palabras reservadas del lenguaje (\code{graph}, \code{kind},
\code{directed}, \code{undirected}, \code{traits}, \code{nodes}, \code{edges},
\code{groups}, \code{constraints}, \code{assert}, los nombres de constraints,
transformaciones y algoritmos, los atributos de nodo, etc.), los operadores
(\code{->}, \code{--}, \code{=}, \code{!=}, \code{>=}, \code{<=}, \code{<},
\code{>}), los delimitadores (\code{:}, \code{\{}, \code{\}}, \code{,},
\code{(}, \code{)}), identificadores (\code{[a-zA-Z\_][a-zA-Z0-9\_]*}),
enteros y comentarios de línea (\code{//}).

\textbf{Análisis sintáctico (Bison).} La gramática (\code{BisonGrammar.y})
define las producciones del lenguaje y las acciones semánticas que construyen el
AST. Es LALR(1) y libre de conflictos.

\textbf{AST (\code{AbstractSyntaxTree.h/c}).} Define los tipos del árbol con un
diseño de unión etiquetada (\emph{tagged union}). El nodo raíz es
\code{Program \{ TopLevelDeclList * decls \}}. Cada \code{TopLevelDecl} es una
unión sobre \code{GraphDecl}, \code{DeriveDecl} y \code{AnalysisDecl}. Se
proveen constructores, destructores y cobertura de \code{\%destructor} en Bison
para evitar \emph{memory leaks}. El AST es el contrato de entrada del backend.

\subsection{Backend}

\subsubsection{Tabla de símbolos}

La tabla de símbolos se implementa mediante \code{IdSet}, un \emph{hash-set} de
\emph{strings} con resolución de colisiones por encadenamiento
(\code{IdSet.h/c}). Ofrece tres operaciones:

\begin{itemize}[noitemsep]
  \item \code{idSetAdd(set, id)} --- inserta una copia del identificador;
        retorna \code{false} si ya existía (detección de duplicados).
  \item \code{idSetContains(set, id)} --- consulta de pertenencia.
  \item \code{idSetForEach(set, visit, ctx)} --- itera sobre todos los
        identificadores (usado para propagar nodos/grupos a grafos derivados).
\end{itemize}

El analizador semántico (\code{SemanticAnalyzer.c}) mantiene las siguientes
tablas, organizadas como una pila de \emph{scopes}:

\begin{center}
\begin{tabular}{@{}lll@{}}
\toprule
\textbf{Tabla} & \textbf{Scope} & \textbf{Contenido} \\
\midrule
\code{graphIds}        & Global       & Grafos declarados y derivados \\
\code{analysisIds}     & Global       & Identificadores de análisis \\
\code{GraphInfo.nodes} & Por grafo    & Nodos declarados \\
\code{GraphInfo.groups}& Por grafo    & Grupos declarados \\
\code{results}         & Por análisis & Resultados de \code{run ... as <id>} \\
\bottomrule
\end{tabular}
\end{center}

Adicionalmente, \code{GraphInfo} almacena el \code{kind} y los \code{traits} de
cada grafo en un registro enlazado, permitiendo validar operaciones
\emph{cross-graph} (derivaciones y análisis que referencian grafos por nombre,
incluyendo grafos derivados).

\subsubsection{Sistema de tipos: \emph{graph-property checking}}

En Nexus no existe un sistema de tipos tradicional con coerción o inferencia.
El ``sistema de tipos'' se reinterpreta como \textbf{validación de propiedades
de grafos}: cada operación impone precondiciones sobre el \code{kind} y los
\code{traits} del grafo objetivo. Se implementaron \textbf{48 validaciones
semánticas}, cada una con un test de rechazo dedicado. La siguiente matriz las
resume agrupadas por categoría y por la función del analizador que las aplica.

\paragraph{Identificadores duplicados.}
\begin{center}
\begin{tabular}{@{}ll@{}}
\toprule
\textbf{Validación} & \textbf{Función} \\
\midrule
Grafo duplicado            & \code{\_checkTopLevel} \\
Análisis duplicado         & \code{\_checkTopLevel} \\
Nodo duplicado             & \code{\_checkGraphScope} \\
Grupo duplicado            & \code{\_checkGraphScope} \\
Resultado duplicado        & \code{\_checkAnalysisScope} \\
Arista paralela duplicada  & \code{\_checkGraphScope} \\
\bottomrule
\end{tabular}
\end{center}

\paragraph{Referencias no declaradas.}
\begin{center}
\begin{tabular}{@{}ll@{}}
\toprule
\textbf{Validación} & \textbf{Función} \\
\midrule
Endpoint de arista       & \code{\_checkGraphScope} \\
Miembro de grupo         & \code{\_checkGraphScope} \\
Grafo fuente de \code{derive} & \code{\_checkTopLevel} \\
Grafo de \code{analysis} & \code{\_checkAnalysisScope} \\
Nodo de \code{constraint}& \code{\_checkConstraints} \\
Grupo de \code{forall}   & \code{\_checkConstraints} \\
Endpoint de \code{run}   & \code{\_checkAnalysisScope} \\
Nodo fuera del subgrafo \code{induced\_subgraph} & \code{\_checkAnalysisScope} \\
Resultado de \code{export} & \code{\_checkAnalysisScope} \\
\bottomrule
\end{tabular}
\end{center}

\paragraph{Traits, operadores y atributos de nodo.}
\begin{center}
\begin{tabular}{@{}ll@{}}
\toprule
\textbf{Validación} & \textbf{Función} \\
\midrule
\code{weight} sin trait \code{weighted}    & \code{\_checkGraphScope} \\
\code{capacity} sin trait \code{capacitated} & \code{\_checkGraphScope} \\
\code{->} en grafo \code{undirected}       & \code{\_checkGraphScope} \\
\code{--} en grafo \code{directed}         & \code{\_checkGraphScope} \\
\code{source}/\code{sink} en \code{undirected} & \code{\_checkGraphScope} \\
Múltiples \code{root}/\code{source}/\code{sink} & \code{\_checkGraphScope} \\
\code{root} con \code{indegree} $\neq 0$   & \code{\_checkGraphScope} \\
\code{source}: \code{indegree}$\neq 0$ u \code{outdegree}$=0$ & \code{\_checkGraphScope} \\
\code{sink}: \code{outdegree}$\neq 0$ o \code{indegree}$=0$ & \code{\_checkGraphScope} \\
\code{terminal} con grado inválido         & \code{\_checkGraphScope} \\
\bottomrule
\end{tabular}
\end{center}

\paragraph{Constraints, transformaciones y algoritmos vs. \code{kind}/\code{traits}.}
\begin{center}
\begin{tabular}{@{}ll@{}}
\toprule
\textbf{Validación} & \textbf{Función} \\
\midrule
\code{connected} en \code{directed}        & \code{\_checkConstraints} \\
\code{strongly\_connected} en \code{undirected} & \code{\_checkConstraints} \\
\code{tree}/\code{binary\_tree} en \code{undirected} & \code{\_checkConstraints} \\
\code{rooted\_at} distinto del nodo \code{root} & \code{\_checkConstraints} \\
\code{acyclic} con ciclo o \emph{self-loop} & \code{\_checkConstraints} \\
\code{transpose}/\code{underlying} en \code{undirected} & \code{\_checkTopLevel} \\
\code{induced\_subgraph} con grupo no declarado & \code{\_checkTopLevel} \\
\code{topological\_sort}/\code{scc} en \code{undirected} & \code{\_checkAlgorithm} \\
\code{topological\_sort} en grafo cíclico  & \code{\_checkAlgorithm} \\
\code{components} en \code{directed}       & \code{\_checkAlgorithm} \\
\code{mst} en \code{directed} o sin \code{weighted} & \code{\_checkAlgorithm} \\
\code{max\_flow} en \code{undirected} o sin \code{capacitated} & \code{\_checkAlgorithm} \\
\code{max\_flow} con endpoints $\neq$ \code{source}/\code{sink} & \code{\_checkAlgorithm} \\
\code{shortest\_path} sin \code{weighted}  & \code{\_checkAlgorithm} \\
\bottomrule
\end{tabular}
\end{center}

\paragraph{Estructurales.}
Grafo sin nodos (\code{\_checkGraphScope}) y análisis sin sentencias
(\code{\_checkAnalysisScope}): aunque la gramática los admite, no son
constructos semánticamente útiles y se rechazan.

\subsubsection{Generación de código}

El generador (\code{Generator.c}) realiza un recorrido del AST dirigido por
sintaxis (\emph{syntax-directed}), emitiendo Python a \emph{stdout}. La función
\code{executeGenerator} recorre la lista de \code{TopLevelDecl} y despacha a:

\begin{itemize}[noitemsep]
  \item \code{\_generateGraph} --- emite la construcción del grafo
        (\code{Graph(directed=..., weighted=..., capacitated=...)}), seguida de
        \code{add\_node}, \code{add\_edge}, \code{add\_group} y las
        \emph{assertions} de las restricciones.
  \item \code{\_generateDerive} --- emite una asignación con la llamada a la
        transformación correspondiente.
  \item \code{\_generateAnalysis} --- emite los \code{run} como asignaciones al
        resultado del algoritmo y los \code{export} como llamadas a
        \code{write\_dot}/\code{write\_json}.
\end{itemize}

El encabezado generado importa los símbolos del \emph{runtime}:

\begin{lstlisting}[language=Python]
from nexus_runtime import Graph
from nexus_runtime import transpose, induced_subgraph, remove_self_loops, underlying
from nexus_runtime import shortest_path, topological_sort, components, scc, mst, max_flow
from nexus_runtime import assert_connected, assert_strongly_connected, assert_acyclic
from nexus_runtime import assert_reachable, assert_tree, assert_binary_tree, assert_forall
from nexus_runtime import write_dot, write_json
\end{lstlisting}

\textbf{Decisión de diseño.} Se optó por emitir un programa Python plano (sin
clases, funciones ni indentación) porque el modelo de ejecución de Nexus es
estrictamente secuencial: declarar, derivar, analizar. Esto simplifica el
generador y hace la salida directamente legible.

\subsubsection{Runtime}

La pregunta de diseño que plantea la consigna --- la relevancia o no de un
\emph{runtime} para el dominio --- se respondió afirmativamente: el dominio de
grafos requiere estructuras de datos y algoritmos no triviales que conviene
delegar a una biblioteca, en lugar de inlinearlos en cada programa generado. El
\emph{runtime} (\code{nexus\_runtime.py}) es una biblioteca Python de
\emph{solo biblioteca estándar} que se distribuye como \emph{asset} junto al
compilador; no se enlaza al compilador, es una dependencia del programa
generado.

\begin{center}
\begin{tabular}{@{}p{0.26\linewidth}p{0.66\linewidth}@{}}
\toprule
\textbf{Componente} & \textbf{Función} \\
\midrule
\code{Graph}   & \code{add\_node}, \code{add\_edge}, \code{add\_group},
                 \code{neighbors}. Representación: lista de adyacencia con
                 diccionarios \code{\{to, weight, capacity\}}. \\
Derivaciones   & \code{transpose}, \code{induced\_subgraph},
                 \code{remove\_self\_loops}, \code{underlying}. \\
Algoritmos     & \code{shortest\_path} (Dijkstra), \code{topological\_sort}
                 (Kahn), \code{components} (BFS), \code{scc} (Tarjan),
                 \code{mst} (Kruskal + Union-Find), \code{max\_flow}
                 (Edmonds-Karp). \\
Assertions     & \code{assert\_connected}, \code{assert\_strongly\_connected},
                 \code{assert\_acyclic}, \code{assert\_reachable},
                 \code{assert\_tree}, \code{assert\_binary\_tree},
                 \code{assert\_forall}. \\
Exportadores   & \code{write\_dot}, \code{write\_result\_dot},
                 \code{write\_json}. \\
\bottomrule
\end{tabular}
\end{center}

\textbf{Validaciones \emph{compile-time} vs. \emph{runtime}.} Las restricciones
del lenguaje se verifican en dos momentos complementarios:
\begin{itemize}[noitemsep]
  \item \emph{Compile-time}: el analizador semántico verifica que la
        restricción sea \emph{aplicable} al tipo de grafo (p.\,ej. \code{tree}
        sólo en \code{directed}). Previene errores de tipo.
  \item \emph{Runtime}: las \emph{assertions} ejecutables verifican que el grafo
        concreto \emph{cumple} la restricción (p.\,ej. que efectivamente es un
        árbol). Previene errores de valor que no pueden conocerse estáticamente.
\end{itemize}

\subsubsection{Pruebas}

La suite de pruebas final comprende:

\begin{center}
\begin{tabular}{@{}lr@{}}
\toprule
\textbf{Suite} & \textbf{Programas} \\
\midrule
Aceptación (frontend + backend)         & 35 \\
Rechazo léxico/sintáctico               & 12 \\
Rechazo semántico                       & 48 \\
Generación de código \emph{end-to-end}  & 24 \\
\bottomrule
\end{tabular}
\end{center}

El \emph{harness} \code{test.sh} ejecuta las suites de aceptación y rechazo
(verificando códigos de salida) y \code{test-codegen.sh} genera Python desde los
programas de aceptación, lo ejecuta y comprueba que se produzcan los artefactos
DOT/JSON esperados. Una \emph{pipeline} de Integración Continua
(\code{.github/workflows/pipeline.yaml}) instala el toolchain, construye y
ejecuta ambos \emph{harness} en cada \emph{push} y \emph{pull request}.

\subsection{Adicionales}

\begin{itemize}[noitemsep]
  \item \textbf{Exportador DOT:} genera archivos compatibles con Graphviz, con
        peso y capacidad como \emph{labels} de arista y atributos de nodo.
  \item \textbf{Exportador JSON:} serializa grafos como
        \code{\{nodes: [...], edges: [...]\}} y resultados de algoritmos como
        estructuras Python.
  \item \textbf{Grafos derivados con validación completa:} al procesar un
        \code{derive}, el analizador sintetiza un \code{GraphInfo} para el grafo
        derivado, copiando nodos y grupos del grafo fuente (vía
        \code{idSetForEach}); \code{underlying} además cambia el \code{kind} a
        \code{undirected}. Esto permite que análisis posteriores referencien
        grafos derivados con la misma validación que los grafos declarados.
  \item \textbf{Metodología TDD:} cada validación semántica se desarrolló
        escribiendo primero el test de rechazo (RED), implementando la regla y
        verificando el paso a verde (GREEN).
\end{itemize}

\subsection{Dificultades Encontradas}

\begin{enumerate}[noitemsep]
  \item \textbf{Resolución de nodos \emph{cross-graph}.} Los grafos derivados
        necesitan conocer los nodos del grafo fuente para validar análisis
        posteriores. Se resolvió con \code{idSetForEach} para copiar los
        conjuntos de nodos y grupos al \code{GraphInfo} del derivado.
  \item \textbf{Logs contaminando \emph{stdout}.} El sistema de \emph{logging}
        del proyecto base envía DEBUG/INFO a \emph{stdout}, donde el generador
        también escribe el código Python. Se resolvió capturando la salida con
        \code{LOGGING\_LEVEL=CRITICAL}.
  \item \textbf{Predicados de grado a \emph{runtime}.} Predicados como
        \code{forall n in g: degree(n) >= K} no pueden verificarse en
        \emph{compile-time} porque requieren la estructura concreta del grafo;
        se delegaron al \emph{runtime} como \emph{assertions} parametrizadas.
  \item \textbf{\emph{Ownership} de memoria bajo Bison.} El AST mantiene el
        \emph{ownership} de todos los \emph{strings} (duplicados con
        \code{strdup} en las acciones semánticas). El analizador y el generador
        usan punteros prestados; los \code{IdSet} duplican sus claves para ser
        independientes del AST.
  \item \textbf{Decisiones gramática vs. semántica.} Algunas restricciones
        podrían haberse forzado en la gramática, pero se optó por una gramática
        permisiva con validación semántica posterior: simplifica la gramática y
        produce mejores mensajes de error.
\end{enumerate}

% ==========================================================================
\section{Futuras Extensiones}

\begin{itemize}[noitemsep]
  \item \textbf{Coerción de tipos:} permitir \code{shortest\_path} en grafos no
        \code{weighted} asumiendo peso unitario.
  \item \textbf{Más algoritmos:} recorridos BFS/DFS, detección de ciclos con
        reporte del ciclo, flujo de costo mínimo.
  \item \textbf{Salida multi-archivo:} emitir cada análisis a un archivo
        separado en vez de un único programa.
  \item \textbf{Lenguaje de \emph{constraints} más rico:} combinaciones lógicas
        (\code{and}, \code{or}, \code{not}) y predicados sobre aristas.
  \item \textbf{\emph{Targets} alternativos:} además de Python, emitir C++,
        Rust o directamente DOT/JSON sin \emph{runtime}.
  \item \textbf{Importación de grafos:} cargar grafos desde archivos DOT o JSON.
  \item \textbf{Subsistema generador/build:} construcción paramétrica de grafos
        con variables locales, condicionales y estructuras iterativas acotadas
        (propuesto ya en Stage I como extensión).
\end{itemize}

% ==========================================================================
\section{Conclusiones}

El compilador Nexus implementa exitosamente las tres etapas del proyecto: un
frontend libre de conflictos que construye un AST completo, un analizador
semántico con 48 validaciones que cubren identificadores, referencias,
\emph{traits}, operadores, atributos de nodo, restricciones, transformaciones,
algoritmos y reglas estructurales, y un generador que produce programas Python
ejecutables.

La decisión de generar Python con un \emph{runtime} de solo biblioteca estándar
resultó acertada: los programas generados se ejecutan sin dependencias externas
(salvo Python 3 y, opcionalmente, Graphviz para visualizar los DOT), lo que
mantiene al compilador libre de dependencias enlazadas y facilita el
\emph{testing} \emph{end-to-end}.

El enfoque TDD fue particularmente efectivo para las validaciones semánticas:
cada regla quedó documentada por un test de rechazo dedicado que describe
exactamente qué programa inválido detecta. La separación entre validaciones de
\emph{compile-time} (aplicabilidad) y \emph{runtime} (cumplimiento) permitió
modelar fielmente el dominio sin sobrecargar ninguna de las dos fases.

% ==========================================================================
\section{Referencias}

\begin{itemize}[noitemsep]
  \item Stage I --- Diseño del lenguaje Nexus: \code{doc/TLA\_Stage\_I.pdf}
        (documento de diseño original).
  \item Flex-Bison-Compiler v2.0.0:
        \url{https://github.com/agustin-golmar/Flex-Bison-Compiler/tree/v2.0.0}
        (proyecto base).
  \item Material de la cátedra: apuntes de Análisis Semántico (v0.2.0), Tabla de
        Símbolos, Sistema de Tipos, Generación de Código (v0.1.0) y Runtime
        (v0.1.0).
\end{itemize}

% ==========================================================================
\section{Bibliografía}

\begin{itemize}[noitemsep]
  \item Aho, A.\,V., Lam, M.\,S., Sethi, R., \& Ullman, J.\,D. (2006).
        \emph{Compilers: Principles, Techniques, and Tools} (2.\textsuperscript{a} ed.).
        Addison-Wesley.
  \item Levine, J. (2009). \emph{flex \& bison}. O'Reilly Media.
  \item Cormen, T.\,H., Leiserson, C.\,E., Rivest, R.\,L., \& Stein, C. (2009).
        \emph{Introduction to Algorithms} (3.\textsuperscript{a} ed.). MIT Press.
\end{itemize}

\end{document}
